\section{Q1: the observations are largely correct}
\label{sec:facts}

This section checks (A1)--(A6). The short answer is that \N\ read the code carefully and reports it accurately, and that the kernel-level export confirms the points on which a text search could mislead.
Table~\ref{tab:facts} at the end of the section collects the individual checks.
Three kinds of evidence are distinguished throughout: \textsf{S}, read by us in the source; \textsf{K}, derived from the kernel export through the dependency graph;
\textsf{G}, established by one of the independent Gemini agents (the letter names the agent; see Section~\ref{sec:agents}).

\subsection{The predicate and its wrapper (A3)}

\excerpt{anygeo}

Two predicates of the same type stand side by side. \lean{geoOf} asks for a positive integer $c$ with \lean{IsGeoOf} and for \lean{HasTrueSurf}; \lean{anyGeometry} returns \lean{True} on every input.
The predicate enters the development through a certificate \lean{OAt} with a field \lean{geo}, and through an inductive closure:

\excerpt{geocl}

The place where the \lean{geo} field is filled with \lean{trivial} is the one \N\ names:

\excerpt{divexprov}

\subsection{The route, and what lies on it (A2, A3)}
\label{sec:route}

\excerpt{main}
\excerpt{zfa}

The lines above are the last two links of the route described in \N, \S7. One level up, the witness \lean{h} is produced by \lean{nonempty_hyps_definite_side_of_closed_bound},
whose proof passes \texttt{geoCl anyGeometry} to three lemmas:

\excerpt{siegel_a}

A text search shows that an identifier is \emph{mentioned}; whether the proof term \emph{uses} it has to be read off the kernel export.

\begin{verdictbox}
\textbf{Kernel-level result (K).} In the dependency graph of the export, exactly \textbf{\NanyGeo} declarations depend on \lean{GlobalTheory.Mem.Top.Bl.anyGeometry}, and all five headline theorems are among them.
One chain is
{\small\lean{riemannHypothesis} $\to$ \lean{rhOfAll} $\to$ \lean{allChi} $\to$ \lean{UpperTheory.FinalTheorems.everyCharacter} $\to$ \lean{…everyCharacter_of_row}
$\to$ \lean{…every_character_of_literature_of_vWin} $\to$ \lean{UpperTheory.Lit.every_character_of_theta_integrals_local_matching}
$\to$ \lean{UpperTheory.Lit.nonempty_hyps_definite_side_of_closed_bound} $\to$ \lean{anyGeometry}.}

The names \lean{plugAll}, \lean{shapeWin}, \leanu{everyE₂}, \lean{linkDIVt_closedK} and \lean{geoOf} do not occur in the export: they are not used by any of the five theorems.
\end{verdictbox}

So (A2) and (A3) hold, including the correction that \N\ itself makes to earlier accounts: the route does not pass through \lean{plugAll} but through its twin \lean{zero_free_all_of_green_v_win}.

One detail of (A3) deserves a remark because \N\ leans on it (\S7, \S10, App.~A). The two arguments \lean{trivial}\ \lean{trivial} in \lean{zero_free_all_of_green_v_win} are not geometric data that has gone missing.
They instantiate two \emph{propositional parameters} of a generic lemma:

\excerpt{shapeof}
\excerpt{topshape}

\lean{forall_goal_of_shape_of_hyps} is stated for arbitrary propositions \lean{V} and \lean{NR}. The shape theorem used on the route is stated with \texttt{V := True} and \texttt{NR := True},
and its proof introduces the two corresponding hypotheses (\lean{hVC} and \lean{nR} in the source) and never refers to them again (\textsf{S}; confirmed independently by agent~H).
The route supplies \lean{trivial} because the statement being applied asks for a proof of \lean{True}.

\subsection{The degenerate model (A4)}
\label{sec:degenerate}

\excerpt{geodegen}

This is an instance of the abstract interface \texttt{DV.Geo F}: the variety is $\mathbb{N}$, and effectivity, nefness, bigness and the like are the proposition \lean{True}.
The file is called \path{Inhabited.lean}; \lean{mdlJunk}, \lean{cxJunk}, \lean{algDegenerate} and \lean{padDegenerate} are defined next to it.
Agent~E checked every row of the table in \N, \S3 (the fields of \lean{geoDegenerate}, of the datum \lean{divXFl} built on it, and of the seed \lean{x0}) against the source and found all of them accurate (\textsf{G-E}).
It also found that the description of \lean{graft} in \N\ is slightly simplified: besides the geometric fields from \lean{divXFl} and the lattice fields from the seed, about fifteen local fields come from a third datum.

\begin{verdictbox}
\textbf{Kernel-level result (K).} All of these objects are in the dependency cone of \lean{riemannHypothesis}. The numbers of declarations that depend on each are:
\lean{geoDegenerate} \NgeoDeg, \lean{algDegenerate} 41, \lean{padDegenerate} 42, \lean{mdlJunk} 44, \lean{cxJunk} 42, \lean{divXFl} 27, \lean{geoOver_stages} 25, the seed \lean{x0} 51, \lean{graft} 28.
The graph shows two further trivial instances of the interface that \N\ does not mention: an empty model \lean{geoExample} (53 dependents) and a one-point model \lean{pointGeo} (47 dependents).
The union of all these sets, together with the dependents of \lean{anyGeometry}, has \textbf{\Nflagged} elements.

The \emph{interface type} \lean{DV.Geo} is a different matter. It is a field of the record \lean{DivX} that is carried through the family data, and \NifaceDeps\ declarations depend on it.
But only \NaccDirect\ declarations refer directly to one of its field accessors, and only \NaccThmProof\ of these are theorems whose proof does so. No declaration eliminates a \lean{DV.Geo} or a \lean{DivX} record by its recursor.
\end{verdictbox}

The picture is therefore the one \N\ describes: a thin layer of declarations is written against an abstract geometry interface, and the release instantiates the interface with trivial models.
Section~\ref{sec:inference} asks what follows from that.

\subsection{The lattice side (A4)}

\excerpt{existsdivt}

The three lattice conditions of the existence statement \lean{DivEx} are discharged as \N, \S4 says: by \texttt{exists\_mjHTie … rfl HEq.rfl}, by \texttt{fun \_ => rfl}, and by \texttt{lamAre\_of\_eq … rfl} (\textsf{S}, \textsf{G-E}).
They can be, because the seed sets \texttt{MjH := mjH Z}, \texttt{logLam := logLamOf Z hH} and \texttt{kapR := fun \_ => kapSharp D}.

Agent~E added a point that \N\ passes over quickly and that matters for (B1). The objects that are copied are not themselves trivial.
\lean{mjLat}, \lean{logLamOf} and \lean{kapSharp} are Hermitian lattices, logarithms of eigenvalues and a spectral constant computed from the arithmetic (``row-side'') data,
and the release proves lemmas with content about them (\lean{mjH_good}, \lean{jM_lat}, \lean{jM_metric}) from hypotheses \lean{Holds} and \lean{Pins} (\textsf{G-E}; all six names are in the export, \textsf{K}).
What holds by \lean{rfl} is that the datum \emph{carries} these lattices, not anything about the lattices.

\subsection{The height (A6)}
\label{sec:height-facts}

\excerpt{hexpl6}
\excerpt{heighteq}

The height slot of the datum is overwritten by \lean{setDiv} with \texttt{p.ht}, which is the formula \lean{hExpl6}; the field \lean{h_eq} then holds by \lean{rfl}. This is (A6), and it is accurate (\textsf{S}, \textsf{G-G}).
Agent~G checked the remaining details of \N, \S7 and reports the following (we re-checked the items marked \textsf{S}):
\begin{itemize}
\item Every caller of \lean{exists_divStage_of_pk} passes the zero function as initial height, as \N\ says (\textsf{G-G}).
\item Of the four sites at which \N\ says \lean{h_eq} is constructed, two (\path{DisjointSplitLevelPlaces.lean}, lines 331 and 501) concern a field of the same name in a different structure,
\texttt{h\_eq : h = hP * k} at line 264 of that file (\textsf{S}); two other sites that propagate the real field are not listed in \N.
\item The construction \lean{linkDIVt_closedK}, which \N\ uses to explain how the height is supplied, is not in the export (\textsf{K}). \N\ notes later in \S7 that the live route uses \lean{closed_linkDIVt_of_closed_bridge_kernel} instead; that theorem is in the export.
\item On the live route the pairing \lean{Bk} is not identically zero: it is tied to the period pairing \lean{Bu} by a proved rationality statement \lean{BuRational}, and the case $\mtt{Bk}\ u \neq 0$ is used downstream (\textsf{G-G}; the lemmas named by the agent are in the export, \textsf{K}).
\end{itemize}
The last sentence of (A6), that the two inequalities about the height are proved from matrix estimates, is also accurate. It is examined in Section~\ref{sec:squeeze} because it decides (B1).

\subsection{The dead extraction lemmas (A5)}
\label{sec:dead}

\N, \S6 reports that \lean{GeoCl} is never opened by case analysis or induction and that six lemmas written to extract geometric content occur only at their own declarations.
A text search cannot quite establish the first claim, since pattern matching need not name the type. The kernel export can.

\begin{verdictbox}
\textbf{Kernel-level result (K).}
\begin{itemize}
\item The recursor \lean{GeoCl.rec} is referenced by no declaration outside the inductive block of \lean{GeoCl} itself; \lean{GeoCl.casesOn} and \lean{GeoCl.recOn} do not occur in the export.
Any case analysis or induction on a \lean{GeoCl} proof compiles to the recursor. \emph{No proof in the cone ever opens a} \lean{GeoCl} \emph{certificate.}
\item The constructor \lean{GeoCl.base} is used once (in \lean{geoCl_of}) and \lean{GeoCl.step} once (in \lean{oAt_setDiv}).
\item The \lean{geo} field of the certificate \lean{OAt} is read by exactly two theorems, \lean{divEx_geoCl} and \lean{oAt_setDiv}, and each passes it straight into a \lean{GeoCl} constructor. \lean{OAt.rec} has no referrer.
By contrast the other field, \lean{lam}, is also read by \lean{GlobalTheory.T.parameter_of_discrete_series_ray_alt}: in the same certificate one field has a consumer and the other has none.
\item \lean{geoOf}, \lean{GeoTiedP.geoTied}, \lean{GeoTiedP.formsTieX}, \lean{degree_value_of_determinant} and \lean{SlopeBoundFamilyStatement} do not occur in the export.
\end{itemize}
\end{verdictbox}

The last item bears on the remark in \N, \S6 that the degree equation is ``never proved''. It occurs only as a hypothesis inside \lean{SlopeBoundFamilyStatement} and one admissibility bundle.
Since \lean{SlopeBoundFamilyStatement} is not in the export, no declaration used by the final theorem assumes that equation. An unproved hypothesis of an unused statement is not a gap. \Fnote

\subsection{Summary of the fact check}

{\small
\renewcommand{\arraystretch}{1.18}
\begin{longtable}{@{}r>{\raggedright\arraybackslash}p{6.5cm}c>{\raggedright\arraybackslash}p{6.3cm}@{}}
\caption{Observations of \N\ and what we found. Evidence: \textsf{S} source, \textsf{K} kernel export, \textsf{G-x} Gemini agent x.}\label{tab:facts}\\
\toprule
\# & Observation in \N\ (section) & Finding & Evidence and remarks \\
\midrule
\endfirsthead
\toprule
\# & Observation in \N\ (section) & Finding & Evidence and remarks \\
\midrule
\endhead
\bottomrule
\endfoot
1 & Closed theorem of Mathlib's \lean{RiemannHypothesis}; no axiom, \lean{sorry}, \lean{native_decide}, \lean{unsafe}, \lean{opaque} (Abstract, \S12) & \OK & \textsf{S}, \textsf{K}: three standard axioms only; statement cone is Mathlib only (Section~\ref{sec:taut}) \\
2 & Route to the final theorem; \lean{plugAll} not on it (\S1, \S7, \S11) & \OK & \textsf{S}, \textsf{K} \\
3 & \texttt{anyGeometry := fun \_ \_ \_ => True}; 118 uses in 24 files (\S2) & \OK & \textsf{S}; \textsf{K}: \NanyGeo\ dependents, final theorem included \\
4 & \lean{divEx_prov} fills the geometric field with \lean{trivial} (\S2) & \OK & \textsf{S}, \textsf{G-E} \\
5 & \lean{divEx_of_ox} ignores three hypotheses named after published results (\S2) & \OK & \textsf{S}: body is \texttt{fun \_ \_ \_ 𝒮 hOx => divEx\_prov 𝒮 hOx}. \textsf{G-E}: at the call sites the three arguments are theorems proved in the release \\
6 & Field table of the degenerate model, of \lean{divXFl} and of the seed (\S3) & \OK & \textsf{G-E}: every row verified; \textsf{K}: all in the cone \\
7 & \lean{graft} takes geometric fields from \lean{divXFl}, lattice fields from the seed (\S3) & \PARTLY & \textsf{G-E}: also about 15 local fields from a third datum \\
8 & Lattice conditions of \lean{DivEx} discharged by \lean{rfl} (\S4) & \OK & \textsf{S}, \textsf{G-E}; the copied lattices are substantive objects with proved lemmas \\
9 & Row-side conditions \lean{Holds}, \lean{InGeo} have content (\S5) & \OK & \N's own statement; \textsf{K}: in the cone\Frow \\
10 & \lean{GeoTiedP.geoTied}, \lean{.formsTieX} never called (\S5) & \OK & \textsf{K}: neither name occurs in the export \\
11 & \lean{GeoCl} never opened (\S6) & \OK & \textsf{K}: \lean{GeoCl.rec} has no referrer outside its own block \\
12 & Six extraction lemmas occur once (\S6) & \OK & \textsf{K}: those checked are absent from the export\Fsix \\
13 & Degree equation only a hypothesis, never proved (\S6) & \OK & \textsf{K}: and the statement containing it is not used by the final theorem \\
14 & Two shape inputs supplied by \lean{trivial} (\S7) & \OK & \textsf{S}, \textsf{G-H}: they instantiate two unused propositional parameters at \lean{True} \\
15 & Height overwritten by \lean{hExpl6}; \lean{h_eq} by \lean{rfl} (\S7) & \OK & \textsf{S}, \textsf{G-G} \\
16 & \lean{h_eq} constructed at four listed sites (\S7) & \PARTLY & \textsf{S}, \textsf{G-G}: two of the four are a different field of the same name \\
17 & Initial height is zero (\S7) & \OK & \textsf{G-G} \\
18 & \lean{WordFacts}: two inequalities, proved by \lean{divClausesR_of} from about forty inputs (\S7) & \OK & \textsf{S}, \textsf{K}, \textsf{G-G}: \NdivInputs\ fields; \lean{WordFacts} has a third conjunct not shown in \N \\
19 & \lean{Bk} is the pairing or else 0 (\S7) & \PARTLY & \textsf{G-G}: on the live route it is tied to the pairing by a proved statement \\
20 & The squeeze ``compares a quantity with itself and holds by reflexivity'' (Abstract, \S8) & \NO & \textsf{S}, \textsf{K}, \textsf{G-G}, \textsf{G-H}: contradicted by row 18; Section~\ref{sec:squeeze} \\
21 & Only the tautology ``RiemannHypothesis $=$ RiemannHypothesis'' is certified (Abstract, \S8) & \NO & \textsf{K}, \textsf{G-H}: Section~\ref{sec:taut} \\
\end{longtable}
}
